% !TEX root = ../libretto.tex
% Sezione: Salmi
\momento{Salmi}
\begin{songs}{titoli}

% Fonte: LibrettoAccordi.pdf (Canti Coro), O23.
% Nell'elenco: Salmo, Pasqua. Messo tra i Salmi (regola dell'utente: i salmi stanno nei Salmi).
\beginsong{Alleluia, rendete grazie (Epistola)}[by={Gen Verde (CD Il mistero pasquale)}]
\beginchorus
Alleluia, alleluia, alleluia,
alleluia, alleluia. \rep{2}
\endchorus
\beginverse
Rendete grazie a Dio, Egli è buono,
eterno e fedele è il suo amore.
Sì, è così: lo dica Israele,
dica che il suo amore è per sempre.
\endverse
\beginverse
La destra del Signore si è innalzata
a compiere grandiose meraviglie.
Non morirò, ma resterò in vita
e annuncerò i prodigi del Signore.
\endverse
\beginverse
La pietra che avevano scartato
è divenuta pietra angolare.
Questo prodigio ha fatto il Signore:
una meraviglia ai nostri occhi.
\endverse
\endsong

% Fonte: LibrettoAccordi.pdf (Canti Coro), D13.
% Nell'elenco: Salmo, Anno A.
\beginsong{Ascoltate oggi}[by={P. Spoladore}]
\beginchorus
Ascoltate oggi la voce del Signore:
non indurite il vostro cuore.
\endchorus
\beginverse
Venite, cantiamo al Signore,
acclamiamo alla roccia della nostra salvezza.
Accostiamoci a lui per rendergli grazie,
a lui acclamiamo con canti di gioia.
\endverse
\beginverse
Entrate: prostrati, adoriamo,
in ginocchio davanti al Signore che ci ha fatti.
È lui il nostro Dio e noi il popolo del suo pascolo,
il gregge che egli conduce.
\endverse
\beginverse
Se ascoltaste oggi la sua voce!
«Non indurite il cuore come a Meriba,
come nel giorno di Massa nel deserto,
dove mi tentarono i vostri padri:
mi misero alla prova pur avendo visto le mie opere».
\endverse
\endsong

% Fonte: LibrettoAccordi.pdf (Canti Coro), D18.
% Nell'elenco: Salmo, Anno B.
\beginsong{Camminerò}[by={P. Spoladore}]
\beginchorus
Camminerò alla presenza del Signore
nella terra dei viventi.
\endchorus
\beginverse
Ho creduto anche quando dicevo:
«Sono troppo infelice».
Agli occhi del Signore è preziosa
la morte dei suoi fedeli.
\endverse
\beginverse
Ti prego, Signore, perché sono tuo servo;
io sono tuo servo, figlio della tua schiava:
tu hai spezzato le mie catene.
A te offrirò un sacrificio di ringraziamento
e invocherò il nome del Signore.
\endverse
\beginverse
Adempirò i miei voti al Signore
davanti a tutto il suo popolo,
negli atri della casa del Signore,
in mezzo a te, Gerusalemme.
\endverse
\endsong

% Fonte: LibrettoAccordi.pdf (Canti Coro), D21.
% Nell'elenco: Salmo, Anno B.
\beginsong{Crea in me, o Dio, un cuore puro}[by={P. Spoladore}]
\beginchorus
Crea in me, o Dio, un cuore puro. \rep{2}
\endchorus
\beginverse
Pietà di me, o Dio, nel tuo amore; nella tua grande misericordia
cancella la mia iniquità.
Lavami tutto dalla mia colpa,
dal mio peccato rendimi puro.
\endverse
\beginverse
Crea in me, o Dio, un cuore puro,
rinnova in me uno spirito saldo.
Non scacciarmi dalla tua presenza
e non privarmi del tuo santo spirito.
\endverse
\beginverse
Rendimi la gioia della tua salvezza,
sostienimi con uno spirito generoso.
Insegnerò ai ribelli le tue vie
e i peccatori a te ritorneranno.
\endverse
\endsong

% Fonte: LibrettoAccordi.pdf (Canti Coro), D1.
% Nell'elenco: Salmo, Pasqua. Messo tra i Salmi (regola dell'utente: i salmi stanno nei Salmi).
\beginsong{Del tuo Spirito, Signore}[by={Gen Verde}]
\beginchorus
Del tuo Spirito, Signore,
è piena la terra, è piena la terra. \rep{2}
\endchorus
\beginverse
Benedici il Signore, anima mia, Signore, Dio, tu sei grande!
Sono immense, splendenti tutte le tue opere e tutte le creature.
\endverse
\beginverse
Se tu togli il tuo soffio muore ogni cosa, e si dissolve nella terra.
Il tuo Spirito scende: tutto si ricrea e tutto si rinnova.
\endverse
\beginverse
La tua gloria, Signore, resti per sempre. Gioisci, Dio, del creato.
Questo semplice canto salga a te Signore, sei tu la nostra gioia.
\endverse
\endsong

% Fonte: LibrettoAccordi.pdf (Canti Coro), D16.
% Nell'elenco: Salmo, Pasqua. Messo tra i Salmi (regola dell'utente: i salmi stanno nei Salmi).
\beginsong{Dio mio, Dio mio}[by={P. Spoladore}]
\beginchorus
Dio mio, Dio mio, perché mi hai abbandonato? \rep{2}
\endchorus
\beginverse
Si fanno beffe di me quelli che mi vedono,
storcono le labbra, scuotono il capo:
«Si rivolga al Signore, lui lo liberi,
lo porti in salvo se davvero lo ama!».
\endverse
\beginverse
Un branco di cani mi circonda,
mi accerchia una banda di malfattori;
hanno scavato le mie mani e i miei piedi.
Posso contare tutte le mie ossa.
\endverse
\beginverse
Si dividono le mie vesti,
sulla mia tunica gettano la sorte.
Ma tu, Signore, non stare lontano,
mia forza, vieni presto in mio aiuto.
\endverse
\beginverse
Annuncerò il tuo nome ai miei fratelli,
ti loderò in mezzo all'assemblea.
Lodate il Signore, voi suoi fedeli,
gli dia gloria tutta la discendenza di Giacobbe,
lo tema tutta la discendenza di Israele.
\endverse
\endsong

% Fonte: LibrettoAccordi.pdf (Canti Coro), D12.
% Nell'elenco: Salmo, Anno A.
\beginsong{Donaci, o Signore}[by={P. Spoladore}]
\beginchorus
Donaci, o Signore, il tuo amore: in te speriamo. \rep{2}
\endchorus
\beginverse
Retta è la parola del Signore e fedele ogni sua opera.
Egli ama la giustizia e il diritto; dell'amore del Signore è piena la terra.
\endverse
\beginverse
Ecco, l'occhio del Signore è su chi lo teme, su chi spera nel suo amore,
per liberarlo dalla morte e nutrirlo in tempo di fame.
\endverse
\beginverse
L'anima nostra attende il Signore: egli è nostro aiuto e nostro scudo.
Su di noi sia il tuo amore, Signore, come da te noi speriamo.
\endverse
\endsong

% Fonte: LibrettoAccordi.pdf (Canti Coro), D26.
% Nell'elenco: Salmo, Anno C.
\beginsong{Grandi cose}[by={P. Spoladore}]
\beginchorus
Grandi cose ha fatto il Signore per noi. \rep{2}
\endchorus
\beginverse
Quando il Signore ristabilì la sorte di Sion,
ci sembrava di sognare.
Allora la nostra bocca si riempì di sorriso,
la nostra lingua di gioia.
\endverse
\beginverse
Allora si diceva tra le genti:
«Il Signore ha fatto grandi cose per loro».
Grandi cose ha fatto il Signore per noi:
eravamo pieni di gioia.
\endverse
\beginverse
Ristabilisci, Signore, la nostra sorte,
come i torrenti del Negheb.
Chi semina nelle lacrime
mieterà nella gioia.
\endverse
\beginverse
Nell'andare, se ne va piangendo,
portando la semente da gettare,
ma nel tornare, viene con gioia,
portando i suoi covoni.
\endverse
\endsong

% Fonte: LibrettoAccordi.pdf (Canti Coro), D2. Ordine delle righe ricostruito e confermato dall'utente.
% Nell'elenco: Salmo 29, Pasqua. Messo tra i Salmi (regola dell'utente: i salmi stanno nei Salmi).
\beginsong{Grida di vittoria}[by={FdV}]
\beginchorus
Grida di vittoria, grida di esultanza! \rep{2}
La destra di Dio si esaltò! \rep{4}
\endchorus
\beginverse
Ti esalto, Signore, mi hai liberato,
non hai permesso che sul mio capo
potessero esultare i miei nemici,
ti rendo grazie, mio Dio.
Signore Dio, a te ho gridato,
Signore Dio, mi hai guarito.
Dagli inferi mi hai fatto risalire,
mi hai dato vita, Signore.
\endverse
\beginverse
Cantate a Dio, o suoi fedeli,
rendete grazie al suo santo nome.
Dura un istante la sua ira,
la sua bontà in eterno.
Alla sera sopraggiunge il pianto
ed al mattino ecco la gioia.
Signore Dio, mi hai salvato,
mi hai liberato dalla morte.
\endverse
\beginverse
Ascolta, Signore, misericordia!
Vieni, Signore, in mio aiuto!
A te io grido, Signore Dio!
Tu dalla morte mi salverai!
Il mio lamento hai mutato in danza,
mi hai messo l'abito della gioia.
Senza fine io canterò,
per sempre, mio Dio, ti loderò.
\endverse
\endsong

% Fonte: LibrettoAccordi.pdf (Canti Coro), D25.
% Nell'elenco: Salmo, Anno C.
\beginsong{Gustate e vedete}[by={P. Spoladore}]
\beginchorus
Gustate e vedete com'è buono il Signore. \rep{2}
\endchorus
\beginverse
Benedirò il Signore in ogni tempo,
sulla mia bocca sempre la sua lode.
Io mi glorio nel Signore:
i poveri ascoltino e si rallegrino.
\endverse
\beginverse
Magnificate con me il Signore,
esaltiamo insieme il suo nome.
Ho cercato il Signore: mi ha risposto
e da ogni mia paura mi ha liberato.
\endverse
\beginverse
Guardate a lui e sarete raggianti,
i vostri volti non dovranno arrossire.
Questo povero grida e il Signore lo ascolta,
lo salva da tutte le sue angosce.
\endverse
\endsong

% Fonte: "4 - Salmi/Il cantico dei redenti.pdf" (melodia con accordi). Isaia 12.
% Nell'elenco: Salmo. Strofe 3 e 4 (tagliate nel foglio) dal testo fornito dall'utente.
\beginsong{Il cantico dei redenti}[by={A. Marani}]
\beginchorus
Il Signore è la mia salvezza
e con lui non temo più,
perché ho nel cuore la certezza:
la salvezza è qui con me.
\endchorus
\beginverse
Ti lodo, Signore, perché
un giorno eri lontano da me,
ora invece sei tornato
e mi hai preso con te.
\endverse
\beginverse
Berrete con gioia alle fonti,
alle fonti della salvezza,
e quel giorno voi direte:
lodate il Signore, invocate il suo nome.
\endverse
\beginverse
Fate conoscere ai popoli
tutto quello che lui ha compiuto,
e ricordino per sempre,
ricordino sempre che il suo nome è grande.
\endverse
\beginverse
Cantate a chi ha fatto grandezze
e sia fatto sapere nel mondo;
grida forte la tua gioia, abitante di Sion,
perché grande con te è il Signore.
\endverse
\endsong

% Fonte: LibrettoAccordi.pdf (Canti Coro), D3.
% Nell'elenco: Salmo, Pasqua. Messo tra i Salmi (regola dell'utente: i salmi stanno nei Salmi).
\beginsong{Il canto del mare}[by={M. Frisina}]
\beginchorus
Cantiamo al Signore, stupenda è la sua vittoria.
Signore è il suo nome. Alleluja!
\endchorus
\beginverse
Voglio cantare in onore del Signore,
perché ha trionfato, alleluja.
Ha gettato in mare cavallo e cavaliere.
Mia forza e mio canto è il Signore.
Il mio Salvatore, il Dio di mio padre,
ed io lo voglio esaltare.
\endverse
\beginverse
Dio è prode in guerra, si chiama Signore.
Travolse nel mare gli eserciti,
i carri d'Egitto sommerse nel Mar Rosso,
abissi profondi li coprono.
La tua destra, Signore, si è innalzata,
la tua potenza è terribile.
\endverse
\beginverse
Si accumularon le acque al tuo soffio,
s'alzaron le onde come un argine.
Si raggelaron gli abissi in fondo al mare.
Chi è come te, o Signore?
Guidasti con forza il popolo redento,
e lo conducesti verso Sion.
\endverse
\endsong

% Fonte: LibrettoAccordi.pdf (Canti Coro), D28.
% Nell'elenco: Salmo.
\beginsong{Il dramma messianico}[by={M. Maneri}]
\beginverse
Perché le genti congiurano,
perché invano cospirano i popoli?
Insorgono i re della terra
congiurano insieme contro il Signore e contro il suo Messia:
«Spezziamo le loro catene,
gettiamo via i loro legami!».
Se ne ride chi abita i cieli,
li schernisce dall'alto il Signore.
Egli parla con ira
e li spaventa con il suo sdegno:
«Io l'ho costituito sovrano sul mio santo monte».
Annunzierò il decreto di Dio. Egli mi ha detto:
«Tu sei mio figlio, oggi io ti ho generato.
Io ti darò in possesso le genti di tutta la terra.
Le spezzerai con scettro di ferro,
come vasi di argilla e le frantumerai».
E ora sovrani siate saggi, istruitevi, giudici della terra;
servite Dio con timore e con tremore esultate.
che non si sdegni e voi perdiate la via.
Improvvisa divampa la sua ira.
Beato chi in lui si rifugia.
\endverse
\endsong

% Fonte: LibrettoAccordi.pdf (Canti Coro), D20.
% Nell'elenco: Salmo, Anno B.
\beginsong{Il ricordo di te, Signore}[by={P. Spoladore}]
\beginchorus
Il ricordo di te, Signore,
è la nostra gioia. \rep{2}
\endchorus
\beginverse
Lungo i fiumi di Babilonia,
là sedevamo e piangevamo
ricordandoci di Sion.
Ai salici di quella terra appendemmo le nostre cetre.
\endverse
\beginverse
Perché là ci chiedevano parole di canto coloro che ci avevano deportato,
allegre canzoni, i nostri oppressori: «Cantateci canti di Sion!».
\endverse
\beginverse
Come cantare i canti del Signore
in terra straniera?
Se mi dimentico di te, Gerusalemme, si dimentichi di me la mia destra.
\endverse
\beginverse
Mi si attacchi la lingua al palato
se lascio cadere il tuo ricordo,
se non innalzo Gerusalemme al di sopra di ogni mia gioia.
\endverse
\endsong

% Fonte: LibrettoAccordi.pdf (Canti Coro), D15.
% Nell'elenco: Salmo, Anno A.
\beginsong{Il Signore è bontà e misericordia}[by={P. Spoladore}]
\beginchorus
Il Signore è bontà e misericordia. Il Signore è bontà e misericordia. \rep{2}
\endchorus
\beginverse
Dal profondo a te grido, o Signore;
Signore, ascolta la mia voce.
Siano i tuoi orecchi attenti
alla voce della mia supplica.
\endverse
\beginverse
Se consideri le colpe, o Signore,
Signore, chi ti può resistere?
Ma con te è il perdono:
così avremo il tuo timore.
\endverse
\beginverse
Io spero, Signore.
Spera l'anima mia, attende la sua parola.
L'anima mia è rivolta al Signore
più che le sentinelle all'aurora.
\endverse
\beginverse
Più che le sentinelle l'aurora.
Israele attenda il Signore,
perché con il Signore è la misericordia e grande è con lui la redenzione.
Egli redimerà Israele da tutte le sue colpe.
\endverse
\endsong

% Fonte: LibrettoAccordi.pdf (Canti Coro), D14.
% Nell'elenco: Salmo, Anno A.
\beginsong{Il Signore è il mio pastore}[by={P. Spoladore}]
\beginchorus
Il Signore è il mio pastore, non manco di nulla. \rep{2}
\endchorus
\beginverse
Il Signore è il mio pastore, non manco di nulla.
Su pascoli erbosi mi fa riposare,
ad acque tranquille mi conduce.
Rinfranca l'anima mia.
\endverse
\beginverse
Mi guida per il giusto cammino a motivo del suo nome.
Anche se vado per una valle oscura
non temo alcun male, perché tu sei con me.
Il tuo bastone e il tuo vincastro mi danno sicurezza.
\endverse
\beginverse
Davanti a me tu prepari una mensa
sotto gli occhi dei miei nemici.
Ungi di olio il tuo capo;
il mio calice trabocca.
\endverse
\beginverse
Sì, bontà e fedeltà mi saranno compagne
tutti i giorni della mia vita,
abiterò ancora nella casa del Signore
per lunghi giorni.
\endverse
\endsong

% Fonte: LibrettoAccordi.pdf (Canti Coro), D23.
% Nell'elenco: Salmo, Anno C.
\beginsong{Il Signore è mia luce}[by={P. Spoladore}]
\beginchorus
Il Signore è mia luce e mia salvezza.
Il Signore è mia luce e mia salvezza. \rep{2}
\endchorus
\beginverse
Il Signore è mia luce e mia salvezza:
di chi avrò timore?
Il Signore è difesa della mia vita:
di chi avrò paura?
\endverse
\beginverse
Ascolta, Signore, la mia voce.
Io grido: abbi pietà di me, rispondimi!
Il mio cuore ripete il tuo invito:
«Cercate il mio volto!» Il tuo volto, Signore, io cerco.
\endverse
\beginverse
Non nascondermi il tuo volto,
non respingere con ira il tuo servo.
Sei tu il mio aiuto, non lasciarmi,
non abbandonarmi, Dio della mia salvezza.
\endverse
\beginverse
Sono certo di contemplare la bontà del Signore
nella terra dei viventi.
Spera nel Signore, sii forte,
si rinsaldi il tuo cuore e spera nel Signore.
\endverse
\endsong

% Fonte: LibrettoAccordi.pdf (Canti Coro), D24.
% Nell'elenco: Salmo, Anno C.
\beginsong{Il Signore ha pietà del suo popolo}[by={P. Spoladore}]
\beginchorus
Il Signore ha pietà del suo popolo.
Il Signore ha pietà del suo popolo. \rep{2}
\endchorus
\beginverse
Benedici il Signore, anima mia,
quanto è in me benedica il suo santo nome.
Benedici il Signore, anima mia,
non dimenticare tutti i suoi benefici.
\endverse
\beginverse
Egli perdona tutte le tue colpe,
guarisce tutte le tue infermità,
salva dalla fossa la tua vita,
ti circonda di bontà e misericordia.
\endverse
\beginverse
Il Signore compie cose giuste,
difende i diritti di tutti gli oppressi.
Ha fatto conoscere a Mosè le sue vie,
le sue opere ai figli d'Israele.
\endverse
\beginverse
Misericordioso e pietoso è il Signore,
lento all'ira e grande nell'amore.
Perché quanto il cielo è alto sulla terra,
così la sua misericordia è potente su quelli che lo temono.
\endverse
\endsong

% Fonte: "18 - Frisian/Lacqua_viva_accordi.pdf" (foglio accordi).
% Nell'elenco: Salmo.
\beginsong{L'acqua viva}[by={M. Frisina}]
\beginchorus
Chi berrà la mia acqua
non avrà più sete in eterno
e quest'acqua sarà per lui
fonte di vita per l'eternità.
\endchorus
\beginverse
Affannati e stanchi,
voi oppressi e poveri venite,
attingete con gioia a Lui
alla sorgente di felicità.
\endverse
\beginverse
Fiumi di acqua viva
sgorgheranno in colui che crederà
nel Signore che dona a noi
l'acqua di vita e di verità.
\endverse
\beginverse
Percuotendo una roccia
dissetasti il popolo in cammino,
fa' che sempre noi camminiam
nel tuo timore e nella fedeltà.
\endverse
\beginverse
Fonte inesauribile,
pace eterna, carità perfetta,
noi a mensa con te sediam,
dolce, immensa, santa Trinità.
\endverse
\beginverse*
Amen.
\endverse
\endsong

% Fonte: LibrettoAccordi.pdf (Canti Coro), D10.
% Nell'elenco: Salmo, Ordinario.
\beginsong{Loda il Signore, anima mia}[by={Gen Rosso}]
\beginchorus
Loda il Signore, anima mia,
loda il Signore, anima mia! \rep{2}
\endchorus
\beginverse
Il Signore rimane fedele per sempre,
Lui fa giustizia agli umiliati.
Agli affamati dà ogni giorno il pane,
e libertà a chi è chiuso in prigione.
\endverse
\beginverse
Il Signore dà luce a chi vive nel buio,
Lui risolleva chi è caduto.
E tutti i giusti colma col suo amore,
ai forestieri dà un tetto sicuro.
\endverse
\beginverse
Il Signore è vicino a chi non ha nessuno,
Lui non sorride a chi è violento.
Dovunque arriva il regno del Signore,
vive in principio, ora e per sempre.
\endverse
\endsong

% Fonte: "99 - Raccolte e libretti accordi/Il mistero pasquale_ gv.pdf" (pag. 61). Testo dal Salmo 30.
% Nell'elenco: Salmo, Settimana Santa (Venerdì Santo). CD: Il mistero pasquale (2011). Messo tra i Salmi (regola dell'utente: i salmi stanno nei Salmi).
\beginsong{Nelle tue mani\soloindice{ (Gen Verde)}}[by={Gen Verde (CD Il mistero pasquale)}]
\beginchorus
Padre, nelle tue mani consegno il mio spirito.
\endchorus
\beginverse
In te, Signore, mi sono rifugiato,
mai sarò deluso;
difendimi per la tua giustizia.
Affido alle tue mani il mio spirito;
tu mi hai riscattato, Signore, Dio fedele.
\endverse
\beginverse
Per i nemici e per i miei vicini
io sono un rifiuto,
l'obbrobrio ed il terrore di tutta la gente,
un corpo rinsecchito, un cadavere;
sono un vaso rotto, un coccio da gettare.
\endverse
\beginverse
Ma io confido in te, o mio Signore:
tu sei il mio Dio
e nelle mani tue sono tutti i miei giorni,
dai miei persecutori mi liberi;
dalle loro mani mi liberi, Signore.
\endverse
\beginverse
Sul tuo servo fa' splendere il tuo volto,
salvami, Signore,
per la tua fedeltà e la misericordia.
Coraggio, il vostro cuore si rianimi;
tutti voi sperate, sperate nel Signore.
\endverse
\endsong

% Fonte: LibrettoAccordi.pdf (Canti Coro), A15.
% Nell'elenco: Salmo, Natale. Messo tra i Salmi (regola dell'utente: i salmi stanno nei Salmi).
\beginsong{Oggi è nato per noi}[by={F. Baggio}]
\beginchorus
Oggi è nato per noi il Salvatore del mondo. \rep{2}
\endchorus
\beginverse
Cantate al Signore un cantico nuovo,
cantate al Signore da tutta la terra.
Cantate al Signore e benedite il suo nome.
\endverse
\beginverse
Cantate ogni giorno la sua salvezza,
in mezzo alle genti narrate la gloria,
a tutti i popoli dite le sue meraviglie.
\endverse
\beginverse
Gioiscano i cieli, esulti la terra,
risuoni il mare e quanto racchiude;
sia festa in campagna come nella foresta.
\endverse
\beginverse
Davanti al Signore che viene dal cielo
per giudicare tutta la terra;
giudicherà con giustizia e nella sua fedeltà.
\endverse
\endsong

% Fonte: LibrettoAccordi.pdf (Canti Coro), D11. Salmo 50 (51).
% Nell'elenco: Salmo, Anno A (I domenica di Quaresima).
\beginsong{Perdonaci, Signore}[by={P. Spoladore}]
\beginchorus
Perdonaci, Signore: abbiamo peccato. \rep{2}
\endchorus
\beginverse
Pietà di me, o Dio, nel tuo amore; nella tua grande misericordia
cancella la mia iniquità.
Lavami tutto dalla mia colpa,
dal mio peccato rendimi puro.
\endverse
\beginverse
Sì, le mie iniquità io le riconosco,
il mio peccato mi sta sempre dinanzi.
Contro di te, contro te solo ho peccato,
quello che è male ai tuoi occhi io l'ho fatto.
\endverse
\beginverse
Crea in me, o Dio, un cuore puro,
rinnova in me uno spirito saldo.
Non scacciarmi dalla tua presenza
e non privarmi del tuo santo spirito.
\endverse
\beginverse
Rendimi la gioia della tua salvezza,
sostienimi con uno spirito generoso.
Signore, apri le mie labbra
e la mia bocca proclami la tua lode.
\endverse
\endsong

% Fonte: testo con accordi fornito dall'utente (accordi tolti). Salmo 15.
% Nell'elenco: "Proteggimi o Dio", Salmo (autore in elenco: Gallotta; l'utente indica Ritinha, Segatori).
\beginsong{Proteggimi, o Dio}[by={Ritinha, Segatori}]
\beginverse
Proteggimi, o Dio: in te io mi rifugio.
Ho detto a lui: «Sei tu il mio Signore,
senza di te non ho bene alcuno».
Nelle tue mani, Signore, è la mia vita!
\endverse
\beginchorus
Tu mi indicherai il sentiero della vita,
gioia piena nella tua presenza,
dolcezza senza fine. \rep{2}
\endchorus
\beginverse
Benedico Dio che m'ha dato consiglio;
anche di notte il cuore m'istruisce.
Innanzi a me sempre il Signore,
sta alla mia destra, non posso vacillare.
\endverse
\beginchorus
Tu mi indicherai il sentiero della vita,
gioia piena nella tua presenza,
dolcezza senza fine. \rep{2}
\endchorus
\beginverse
Mia eredità, mio calice è il Signore,
per me la sorte è su luoghi deliziosi.
Lieto e sereno è il cuore mio,
luce e speranza ai miei passi tu darai.
\endverse
\beginchorus
Tu mi indicherai il sentiero della vita,
gioia piena nella tua presenza,
dolcezza senza fine. \rep{2}
\endchorus
\endsong

% Fonte: LibrettoAccordi.pdf (Canti Coro), D22. Salmo 90 (91).
% Nell'elenco: Salmo, Anno C (I domenica di Quaresima).
\beginsong{Resta con noi}[by={P. Spoladore}]
\beginchorus
Resta con noi, Signore, nell'ora della prova, Signore.
\endchorus
\beginverse
Chi abita al riparo dell'Altissimo
passerà la notte all'ombra dell'Onnipotente.
Io dico al Signore: «Mio rifugio e mia fortezza,
mio Dio in cui confido».
\endverse
\beginverse
Non ti potrà colpire la sventura,
nessun colpo cadrà sulla tua tenda.
Egli per te darà ordine ai suoi angeli
di custodirti in tutte le tue vie.
\endverse
\beginverse
Sulle mani essi ti porteranno,
perché il tuo piede non inciampi nella pietra.
Calpesterai leoni e vipere,
schiaccerai leoncelli e draghi.
\endverse
\beginverse
«Lo libererò, perché a me si è legato,
lo porrò al sicuro, perché ha conosciuto il mio nome.
Mi invocherà e io gli darò risposta; nell'angoscia io sarò con lui,
lo libererò e lo renderò glorioso».
\endverse
\endsong

% Fonte: LibrettoAccordi.pdf (Canti Coro), D9.
% Nell'elenco: Salmo.
\beginsong{Salmo 50}[by={RnS}]
\beginverse
Pietà di me, o Dio, secondo la tua misericordia;
nel tuo grande amore cancella il mio peccato.
Lavami dalle mie colpe, mondami dal mio peccato.
Riconosco la mia colpa, il mio peccato mi sta davanti.
\endverse
\beginchorus
Nel tuo amore fa' grazia a Sion, rialza le mura di Gerusalemme. \rep{2}
\endchorus
\beginverse
Crea in me, o Dio, un cuore puro;
rinnova in me uno spirito saldo.
Non allontanarmi dalla tua presenza,
e non privarmi del tuo Santo Spirito.
\endverse
\beginverse
Purificami con issopo ed io sarò mondato;
lavami e sarò più bianco della neve.
Fammi sentire gioia e letizia,
esulteranno le ossa che hai spezzato.
\endverse
\endsong

% Fonte: LibrettoAccordi.pdf (Canti Coro), D19. Salmo 18 (19).
% Nell'elenco: Salmo, Anno B (III domenica di Quaresima).
\beginsong{Signore, tu hai parole di vita eterna}[by={P. Spoladore}]
\beginchorus
Signore, tu hai parole di vita eterna.
Signore, tu hai parole di vita eterna.
\endchorus
\beginverse
La legge del Signore è perfetta,
rinfranca l'anima;
la testimonianza del Signore è stabile,
rende saggio il semplice.
\endverse
\beginverse
I precetti del Signore sono retti,
fanno gioire il cuore;
il comando del Signore è limpido,
illumina gli occhi.
\endverse
\beginverse
Il timore del Signore è puro,
rimane per sempre;
i giudizi del Signore sono fedeli,
sono tutti giusti.
\endverse
\beginverse
Più preziosi dell'oro,
di molto oro fino,
più dolci del miele
e di un favo stillante.
\endverse
\endsong

% Fonte: LibrettoAccordi.pdf (Canti Coro), D27. Salmo 29 (30).
% Nell'elenco: Salmo.
\beginsong{Ti esalterò}[by={P. Spoladore}]
\beginchorus
Ti esalterò, Signore! Ti esalterò, Signore! \rep{2}
Perché mi hai risollevato. \rep{4}
\endchorus
\beginverse
Ti esalterò, Signore, perché mi hai risollevato,
non hai permesso ai miei nemici di gioire su di me.
Signore mio Dio, a te ho gridato e mi hai guarito.
Signore, hai fatto risalire la mia vita dagli inferi.
Mi hai fatto rivivere perché non scendessi nella fossa.
\endverse
\beginverse
Cantate inni al Signore, o suoi fedeli,
della sua santità celebrate il ricordo.
Perché la sua collera dura un istante,
la sua bontà per tutta la vita.
Alla sera è ospite il pianto, al mattino la gioia, la gioia.
\endverse
\beginverse
Ascolta, Signore, abbi pietà di me,
Signore, vieni in mio aiuto.
Hai mutato il mio lamento in danza,
Signore mio Dio, ti renderò grazie.
Signore mio Dio, ti renderò grazie per sempre.
\endverse
\endsong

% Fonte: LibrettoAccordi.pdf (Canti Coro), D17. Salmo 24 (25).
% Nell'elenco: Salmo, Anno B (I domenica di Quaresima).
\beginsong{Tutti i sentieri}[by={P. Spoladore}]
\beginchorus
Tutti i sentieri del Signore
sono amore e fedeltà, sono amore e fedeltà.
\endchorus
\beginverse
Fammi conoscere, Signore, le tue vie,
insegnami i tuoi sentieri.
Guidami nella tua fedeltà e istruiscimi,
perché sei tu il Dio della mia salvezza.
\endverse
\beginverse
Ricordati, Signore, della tua misericordia
e del tuo amore, che è da sempre.
Ricordati di me nella tua misericordia,
per la tua bontà, Signore.
\endverse
\beginverse
Buono e retto è il Signore,
indica ai peccatori la via giusta;
guida i poveri secondo giustizia,
insegna ai poveri la sua via.
\endverse
\endsong

% Fonte: LibrettoAccordi.pdf (Canti Coro), A20 (li' tra i canti d'ingresso).
% Nell'elenco: Salmo, Ordinario. CD: Vita nuova con te (1995).
\beginsong{Un cuore nuovo}[by={F. Buttazzo (CD Vita nuova con te)}]
\beginchorus
Ti darò un cuore nuovo, popolo mio,
il mio Spirito effonderò in te.
Toglierò da te il cuore di pietra,
un cuore di carne ti darò, popolo mio.
\endchorus
\beginverse
Da tutte le nazioni vi radunerò,
vi mostrerò la strada della vita.
E vivrà chi la seguirà.
\endverse
\beginverse
Vi aspergerò con acqua, e puri vi farò,
dagli idoli sarete liberati.
Questa è la mia libertà.
\endverse
\beginverse
Mio popolo sarete, le genti lo vedranno,
abiterete dentro la mia casa,
e vedrete il mio volto.
\endverse
\endsong

\end{songs}
